\documentclass[10pt,a4paper]{article}
\usepackage{packages} % Importing all the packages

\begin{document}

\maketitle

\hrule
\tableofcontents
\vspace{0.5cm}\hrule\vspace{0.5cm}


\section{Something to remember}

Please remember that smaller figures will require larger fonts in order to be readable.

The command \texttt{texttt} allows you to write in a coding-format.

Command \texttt{clearpage} ends the current page and causes all figures and tables that have so far appeared in the input to be printed.

When you want import something in a different folder you have to use \texttt{../src/example.cpp} . To report the output of the code, use can use the \texttt{verbatim} environment. But, if you want insert code use the \texttt{lstlisting} environment included with the \texttt{listing} package included at the beginning of this document. An example is in the template. For instance:
\begin{verbatim}
  \begin{lstlisting}[language=C++]
      double i;
      for (i = ... ){
         ...
      }
      return;
  \end{lstlisting}

\end{verbatim}

If your code is long or you need to include more than one file, then a more efficient strategy is to use the \texttt{\textbackslash lstinputlisting[language=C++]\{fname.cpp\}} command. This is show in the section:


\newpage
\appendix
\section{Inserting Code}
%\section{\texttt{ParticleEM.cpp}}
\lstinputlisting[language=C++]{../src/ParticleEM.cpp}

\end{document}

